We propose the Multiple Oracle algorithm (\ref{alg:mo}) as an iterative strategy-generation technique for (approximately) solving any continuous game~$\calG=\langle N,(X_i)_{i\in N},(u_i)_{i\in N}\rangle$. We recall that the~\emph{best response set} of player $i\in N$ with respect to a mixed strategy profile $\bmu_{-i}\in \bM_{-i}$ is
\[
\beta_i(\bmu_{-i}) =  \argmax_{x\in X_i} U_i(x,\bmu_{-i}).
\]
The set $\beta_i(\bmu_{-i})$ is always nonempty by Proposition \ref{prop:switch}. We assume that every player uses an oracle to recover at least one best response strategy, which means that the player is able to solve the corresponding optimization problem to global optimality. In \cref{sec:custom.oracles} we will see the instances of games for which this is possible.

\begin{algorithm}
	\renewcommand{\thealgorithm}{MO}
	\caption{The Multiple Oracle algorithm}
	\label[mo]{alg:mo}
	\begin{algorithmic}
		\Function{iterate}{$\calG$, $\bX$}
			\State Find an equilibrium $\bmu$ of the subgame $\langle N,\bX,(u_i)_{i\in N}\rangle$
			\State $\forall i \in N:\;$ Pick a best response $x_i\in \beta_i(\bmu_{-i})$
			\State $\forall i \in N:\;$ Extend strategy set $Y_i \gets X_i \cup \{x_i\}$
			\State \Return $\bY, \bx, \bmu$
		\EndFunction
		\Function{multiple-oracle}{$\calG$, $\bX$, $\eps$}
			\Repeat
				\State $\bX, \bx, \bmu \gets \text{iterate}(\calG, \bX)$
			\Until{$\bU(\bx,\bmu_{-i})-\bU(\bmu) \leq \beps$}
			\State \Return $\bmu$
		\EndFunction
	\end{algorithmic}
\end{algorithm}

\cref{alg:mo} proceeds as follows: In every iteration $j$, finite strategy sets $X_i^j$ are constructed for each player $i\in N$ and the corresponding finite subgame of $\calG$ is solved. Let $\bmu^j$ be its equilibrium. Then, an arbitrary best response strategy $x_i^{j+1}$ with respect to $\bmu_{-i}^{j}$ is added to each strategy set~$X_i^j$. Those steps are repeated until
\begin{equation}\label{def:diff}
	U_i(x_i^{j+1},\bmu_{-i}^j)-U_i(\bmu^j)\leq \eps \qquad \forall i\in N.
\end{equation}

First we discuss basic properties of the algorithm.  At each step~$j$, we have the inequality
\begin{equation}\label{ineq:upper}
	\bU(\bx^{j+1},\bmu^j)\ge \bU(\bmu^j).
\end{equation}
Indeed, for each player $i\in N$, we get
\[
	U_i(x_i^{j+1},\bmu_{-i}^j) = \max_{x \in X_i} U_i(x,\bmu^j_{-i}) \ge \max_{x \in X_i^j} U_i(x,\bmu^j_{-i}) = U_i(\bmu^j).
\]

We note that \eqref{def:diff} is necessarily different from the stopping condition of DO. Namely, the~latter condition \cite{McMahan03}, tailored to two-player zero-sum games, is
\begin{equation} \label{stoppingDO}
	U_1(x_1^{j+1},\mu_2^j) - U_1(x_2^{j+1},\mu_1^j) \le \eps.
\end{equation}
When $\calG$ is a two-player zero-sum continuous game, it follows immediately from \eqref{ineq:upper} that \eqref{stoppingDO} implies \eqref{def:diff}.

We prove correctness of \cref{alg:mo} --- the eventual output $\bmu^j$ is an $\eps$-equilibrium of the original game $\calG$.
\begin{lemma}\label{lem:algstop}
	The strategy profile $\bmu^j$ is an $\eps$-equilibrium of $\calG$, whenever \cref{alg:mo} terminates at step $j$.
\end{lemma}
\begin{proof}
	Let $\bx\in \bX$. We get
	\[
		\bU(\bx,\bmu^j)-\bU(\bmu^j) = \bU(\bx,\bmu^j)-\bU(\bx^{j+1},\bmu^j) +  \bU(\bx^{j+1},\bmu^j)-\bU(\bmu^j).
	\]
	Then $\bU(\bx,\bmu^j)-\bU(\bx^{j+1},\bmu^j)\leq \mathbf{0}$, since $\bx^{j+1}$ is the profile of best response strategies. Consequently, we obtain
	\[
		\bU(\bx,\bmu^j)-\bU(\bmu^j)  \le  \bU(\bx^{j+1},\bmu^j)-\bU(\bmu^j) \le \beps,
	\]
	where the last inequality is just the terminating condition \eqref{def:diff}. Therefore, $\bmu^j$ is an $\eps$-equilibrium of  $\calG$.
\end{proof}

If the set $\bX^j=X^j_1\times \dots \times X^j_n$ cannot be further inflated, \cref{alg:mo} terminates.
\begin{lemma}\label{lemma:do_stop}
	Assume that $\bX^j=\bX^{j+1}$ at step $j$ of \cref{alg:mo}. Then
	\[
		\bU(\bx^{j+1},\bmu^j)=\bU(\bmu^j).
	\]
\end{lemma}
\begin{proof}
	The condition $\bX^j=\bX^{j+1}$ implies $\bx^{j+1}\in \bX^{j}$. For each player $i\in N$,
	\[
		U_i(x_i^{j+1},\bmu^j_{-i})=\max\limits_{x\in X_i^j}U_i(x,\bmu^j_{-i}) =U_i(\bmu^j).
	\]
\end{proof}

The original Double Oracle algorithm for two-player zero-sum finite games \cite{McMahan03} makes it possible to find exact equilibria. The following proposition generalizes this result to the setting of multiplayer general-sum finite games.
\begin{proposition}\label{pro:Mahan}
	If $\calG$ is a finite game and $\eps=0$, then \cref{alg:mo} recovers an equilibrium of $\calG$ in finitely-many steps.
\end{proposition}
\begin{proof}
	Let $\calG$ be finite and $\eps=0$. Since each $X_i$ is finite, there exists an iteration $j$  in which $\bX^{j+1}=\bX^j$. Then the terminating condition of \cref{alg:mo} is satisfied with $\eps=0$ (Lemma \ref{lemma:do_stop}) and $\bmu^j$ is an equilibrium of $\calG$ (Lemma~\ref{lem:algstop}).
\end{proof}

Theorems \ref{thm:convergence1} and \ref{thm:convergence2} extend the convergence theorem from \cite{Adam21} and provides an alternative proof of Glicksberg's theorem \cite{Glicksberg52} about existence of equilibria in continuous games. Indeed, the sequence $\bmu^1,\bmu^2,\dots$ has at least one accumulation point by Proposition \ref{prop:convergence}. Further, the consequence of Theorem \ref{thm:convergence2} is that a finitely-supported $\eps$-equilibrium of $\calG$ can always  be found in finitely-many steps, although game $\calG$ may have no finitely-supported equilibria.

\begin{theorem}\label{thm:convergence1}
	Let $\calG$ be a continuous game and $\eps=0$. If \cref{alg:mo} stops at step $j$, then $\bmu^j$ is an~equilibrium of $\calG$. Otherwise, any accumulation point of $\bmu^1,\bmu^2,\dots$ is an equilibrium of $\calG$.
\end{theorem}
\begin{proof}
	If \cref{alg:mo} terminates at step $j$, then Lemma \ref{lem:algstop} implies that $\bmu^j$ is an equilibrium of $\calG$. In the opposite case, the algorithm generates a sequence of mixed strategy profiles $\bmu^1, \bmu^2, \dots$ Consider any weakly convergent subsequence of this sequence --- at least one such subsequence exists by Proposition~\ref{prop:convergence}. Without loss of generality, such a subsequence will be denoted by the same indices as the original sequence. Therefore, for each player $i\in N$, there exists some $\mu^\star_i\in \calM(X_i)$, such that the sequence $\mu_i^1,\mu_i^2,\dots$ weakly converges to $\mu_i^\star$. We need to show that  $\bmu^\star=(\mu^\star_1,\dots,\mu^\star_n)\in \bM$ is an equilibrium of $\calG$.

	For all $i\in N$, define
	\[
		Y_i = \bigcup_{j=1}^{\infty} X^j_i.
	\]
	First, assume that $x\in Y_i$. Then there exists $j_0$ with $x\in X_i^j$ for each $j\ge j_0$. Hence the inequality $U_i(\bmu^j) \geq U_i(x,\bmu_{-i}^j)$, for each $j\ge j_0$, since $\bmu^j$ is an equilibrium of the corresponding finite subgame. Therefore,
	\[
		U_i(\bmu^\star) = \lim_{j}U_i(\bmu^j) \geq \lim_{j} U_i(x,\bmu_{-i}^j) = U_i(x,\bmu_{-i}^\star).
	\]
	Further, by continuity of $U_i$,
	\begin{equation}\label{ineq:NExcl}
		U_i(\bmu^\star) \geq  U_i(x,\bmu_{-i}^\star) \quad \text{for all } x\in\cl(Y_i).
	\end{equation}
	Now, consider an arbitrary $x\in X_i\setminus \cl(Y_i)$. The definition of $x_{i}^{j+1}$ yields $U_i(x_{i}^{j+1},\bmu^j_{-i}) \geq U_i(x,\bmu^j_{-i})$ for each $j$. This implies, by continuity,
	\begin{equation}\label{eq:conv_aux3}
		\lim_{j} U_i(x_{i}^{j+1},\bmu_{-i}^\star) \geq \lim_{j} U_i(x,\bmu^j_{-i}) =  U_i(x,\bmu_{-i}^\star).
	\end{equation}
	Since $x_{i}^{j+1}\in X_i^{j+1}$, compactness of $X_i$ provides a convergent subsequence (denoted by the same indices) such that $ x'= \lim\limits_j x_i^{j} \in \cl(Y_i)$. Then (\ref{ineq:NExcl}) gives
	\begin{equation}\label{ineq:NEx}
		U_i(\bmu^\star)\ge U_i(x',\bmu_{-i}^\star) =\lim_{j} U_i(x_{i}^{j+1},\bmu_{-i}^\star).
	\end{equation}
	Combining (\ref{eq:conv_aux3}) and (\ref{ineq:NEx}) shows that $U_i(\bmu^\star) \geq  U_i(x,\bmu_{-i}^\star)$.
\end{proof}

\begin{theorem}\label{thm:convergence2}
	Let $\calG$ be a continuous game and $\eps>0$. Then \cref{alg:mo} terminates at some step $j$ and $\bmu^j$ is an $\eps$-equilibrium of~$\calG$.
\end{theorem}
\begin{proof}
	If \cref{alg:mo} terminates at step~$j$ with a positive $\eps$, then Lemma \ref{lem:algstop} implies that $\bmu^j$ is an $\eps$-equilibrium of $\calG$. Otherwise \cref{alg:mo} produces a sequence $\bmu^1,\bmu^2,\dots$ and we can repeat the analysis as in Item 1 for  convergent subsequences of $(\bmu^j)$ and $(\bx^j)$, which are denoted by the same indices. Define $x'= \lim_j x_i^j.$ Then
	\begin{equation}\label{ineq:aux1}
	U_i(\bmu^\star)\ge U_i(x',\bmu_{-i}^\star)= \lim_{j} U_i(x_{i}^{j+1},\bmu_{-i}^j).
	\end{equation}
	At every step $j$ we have $U_i(x_i^{j+1},\bmu^j_{-i})\geq U_i(\bmu^j)$ for each $i\in N$ by \eqref{ineq:upper}. For this reason $\lim_{j} U_i(x_{i}^{j+1},\bmu_{-i}^j) \geq  U_i(\bmu^\star).$ Putting together the last inequality with (\ref{ineq:aux1}), we get
	\begin{equation}\label{eq:limiting}
	\lim_{j} (U_i(x_{i}^{j+1},\bmu_{-i}^j)-U_i(\bmu^j)) = 0
	\end{equation}
	for each $i\in N$. This equality means that \cref{alg:mo} stops at some step $j$ and $\bmu^j$ is an $\eps$-equilibrium by Lemma \ref{lem:algstop}.
\end{proof}

\cref{alg:mo} generates the sequence of equilibria $\bmu^1,\bmu^2,\dots$ in increasingly larger subgames of $\calG$. The sequence itself may fail to converge weakly in $\bM$ even for a~two-player zero-sum continuous game; see Example 1 from \cite{Adam21}. In fact, Theorems \ref{thm:convergence1} and \ref{thm:convergence2} guarantee convergence only to an accumulation point. We recall that this is a~typical feature of some globally convergent methods not only in infinite-dimensional spaces \cite[Theorem 2.2]{Hinze08}, but also in Euclidean spaces. For example, a gradient method generates the sequence such that only its accumulation points are guaranteed to be stationary points \cite[Proposition 1.2.1]{Bertsekas16}. One necessary condition for the weak convergence of $\bmu^1,\bmu^2,\dots$ is easy to formulate using the stopping condition of \cref{alg:mo}.
\begin{proposition}
  	If the sequence $\bmu^1,\bmu^2,\dots$ generated by \cref{alg:mo} converges weakly to an equilibrium $\bmu^\star$, then $\lim_j \bigl(U_i(x_i^{j+1},\bmu_{-i}^j)-U_i(\bmu^j)\bigr)=0,$ for each $i\in N.$
\end{proposition}
\begin{proof}
	Using \eqref{ineq:upper} and the triangle inequality,
	\begin{align*}
		U_i(x_i^{j+1},\bmu_{-i}^j)-U_i(\bmu^j)&=\abs{U_i(x_i^{j+1},\bmu_{-i}^j)-U_i(\bmu^j)} \\
		&\leq \abs{U_i(x_i^{j+1},\bmu_{-i}^j)-U_i(\bmu^\star)} + \abs{U_i(\bmu^\star)-U_i(\bmu^j)}.
	\end{align*}
	As $j\to \infty$, the first summand goes to zero by \eqref{eq:limiting} and the second by the assumption. Hence the conclusion.
\end{proof}

In our numerical experiments (see \cref{sec:custom.oracles}), we compute the difference
\begin{equation}\label{crit1}
	U_i(x_i^{j+1},\bmu_{-i}^j)-U_i(\bmu^j)
\end{equation}
at each step $j$ and check if such differences are diminishing with $j$ increasing. This provides a simple heuristics to detect the quality of approximation and convergence. Another option is to calculate the Wasserstein distance
\begin{equation}\label{crit2}
	d_W(\bmu^j,\bmu^{j+1}),
\end{equation}
which can be done with a linear program \eqref{def:LPobj}. If the sequence $\bmu^1,\bmu^2,\dots$ converges weakly, then $d_W(\bmu^j,\bmu^{j+1}) \to 0$. We include the values \eqref{crit2} in the results of some numerical experiments and observe that they are decreasing to zero quickly. Figure \ref{fig:oracle.zerosum.polynomial} shows, that neither \eqref{crit1} nor \eqref{crit2} are monotone sequences.




\cref{alg:mo} is a meta-algorithm parameterized by
\begin{enumerate}
	\item the algorithm for computing equilibria of subgames (the \emph{master problem}) and
	\item the optimization method for computing the best response (the \emph{oracles}).
\end{enumerate}
We detail this setup for each example in the next section.

The choice of computational methods should reflect the properties of a~continuous game, since the efficiency of methods for solving the master problem and subproblem is the decisive factor for the overall performance and precision of \cref{alg:mo}. For example, polymatrix games are solvable in polynomial time \cite{Cai16}, whereas finding even an approximate Nash equilibrium of a finite general-sum game is a very hard problem \cite{Daskalakis09}. As for the solution of the oracles, the best response computation can be based on global solvers for special classes of utility functions.
